\chapter{The Dark Sector: Dark Matter, Dark Energy, and the CMB}
\label{chap:darksector}

\section{Overview}

Three of the most puzzling features of observational cosmology---dark matter,
dark energy, and the CMB---receive unified interpretations in the \SCT
framework. All three are signatures of the supraverse touching our universe
through the Cartasis membrane.

\section{Dark matter as parent's gravitational influence}

\subsection{The mechanism}

The Cartasis membrane persists throughout the parent's lifetime, with new matter
continuously falling through it from the parent's exterior. The matter content on
the parent's side of the membrane exerts gravitational influence on our universe
through the membrane, even though it does not enter our 3D space as ordinary
matter. This gravitational influence manifests as:

\begin{itemize}
  \item halo-like distributions around galaxies (where mass concentrated at
    bounce time);
  \item large-scale structure correlations (parent's mass distribution projected
    through the bounce);
  \item anomalously high rotation velocities of galaxies (additional
    gravitational binding);
  \item gravitational lensing without visible mass.
\end{itemize}

This matches the phenomenology of dark matter.

\paragraph{A rate behind it.} The parent is not a static reservoir---it is a
growing black hole. Holes have negative heat capacity, so above
$M_{\mathrm{crit}} = \hbar c^3/(8\pi G k_B T_{\mathrm{bath}})$ (about
$5\times10^{22}\,\mathrm{kg}$ in today's CMB) they accrete faster than they
evaporate. Every astrophysical parent is therefore growing, and the
parity-violating Cartasis filter (Chapter~\ref{chap:chirality}) keeps acting on
what it accretes, so the injection through the membrane is both \emph{ongoing}
and \emph{of a definite, filtered chirality}. Dark matter is thus the live
gravitational signature of the parent's continuing biased accretion, not a frozen
relic---which is what gives the dark sector a time-dependence (next).

\paragraph{The meal is lumpy: a large-scale structure imprint.} What the parent
feeds in is not smooth. A growing black hole eats whatever its environment
supplies---infalling streams, merging clumps, the cosmic web of \emph{its}
universe---so the matter projected through the membrane carries the parent's
inhomogeneity. For an OG universe this does not arise (the void has no structured
meal), but for any non-OG universe---ours included
(Section~\ref{sec:notogu})---the dark-matter background should retain large-scale
features inherited from the parent's accretion geometry, on scales unrelated to
our own internal structure-formation history. This is a concrete reading of
otherwise puzzling ultra-large structures: the Giant Arc and the
$\sim\!400\,\mathrm{Mpc}$ ``Big Ring'' of \citet{lopez2024bigring}, which sit
uncomfortably with the cosmological principle in $\Lambda$CDM, are natural as frozen
imprints of the lumpiness of the parent's meal. The framework's prediction is
specific: such structures should correlate with the dark-matter background rather
than the luminous one, should appear on scales above the internal homogeneity
scale, and should not require any internal mechanism to assemble in the available
time---the organization is inherited, not built. This is speculative and
needs a quantitative correlation-function test (Chapter~\ref{chap:tests}), but it
is falsifiable and it is the kind of thing $\Lambda$CDM does not naturally produce.

\subsection{Specifically: filter-rejected antimatter}

If the bounce has any filter character (chirality-selective transmission,
Chapter~\ref{chap:chirality}), then the rejected chirality content remains at the
Cartasis membrane. From our perspective the rejected content is gravitationally
present (it has mass--energy), does not appear in our 3D space (it is
geometrically at the membrane, not in our interior), does not interact
electromagnetically with our matter (no normal photon exchange across the
membrane), and does not annihilate with our matter (geometric separation
prevents direct contact).

\textbf{This is the standard phenomenology of dark matter.} The ``missing
antimatter'' of standard cosmology and the ``dark matter'' of observational
cosmology may be the same phenomenon.

\subsection{Quantitative prediction}

If filter efficiency is $f$ (fraction of input matter that passes), then the
dark-matter-to-baryon ratio is approximately $(1-f)/f$. Our observed ratio of
$\sim 5\!:\!1$ implies $f \sim 1/6$. This is a specific number that should be
calculable from Einstein--Cartan parameters. If the calculation gives
$f \approx 1/6$ from first principles, that is a strong validation; if it gives
something very different, the framework needs revision.

\subsection{Predictions distinguishing this from particle dark matter}

\begin{enumerate}
  \item \textbf{Dark matter has chirality structure} (specifically, opposite to
    our matter). Probing dark-matter weak-interaction couplings (if any) should
    reveal antiparticle-like behavior, not particle-like.
  \item \textbf{Dark-matter distribution should track the parent's matter
    structure}, not internal baby dynamics. This might explain anomalies in
    dark-matter halo profiles vs.\ CDM predictions.
  \item \textbf{Dark-matter density should evolve with cosmic time} as the
    parent's accretion history changes. Tensions in dark matter inferred at
    different redshifts (e.g.\ the $\sigma_8$ tension) might be evidence for this.
  \item \textbf{No dark-matter particles exist as fundamental constituents.}
    Direct-detection experiments will continue to find nothing. The dark matter
    is geometric, not particulate.
  \item \textbf{Gravitational waves should propagate normally} (they are features
    of our spacetime), but might show subtle dispersion effects from
    membrane-mediated coupling at very low frequencies.
\end{enumerate}

\section{Dark energy as parent accretion dynamics}

\subsection{The mechanism}

The parent black hole, embedded in its own universe, accretes from its
surroundings. As parent mass grows, the Cartasis membrane geometry changes---
specifically, the membrane's ``size'' relative to the baby interior grows. From
inside the baby, this manifests as additional expansion pressure beyond what
internal baby dynamics would predict. The acceleration of cosmological expansion
(attributed to dark energy in standard cosmology) is actually the baby's response
to parent growth.

This makes dark energy a \emph{rate}: to leading order its strength tracks the
time-derivative of the parent's accretion history. A parent still feeding from a
rich environment sustains or strengthens the internal dark energy; one whose
surroundings are depleting, or which is turning toward Hawking-dominated mass
loss, produces weakening dark energy. The observed evolution of dark energy is
therefore a direct, if model-dependent, readout of our parent's growth---the
natural place to confront the DESI~\citep{desi2024} $w_0$--$w_a$ hints without
forcing a sign.

\subsection{Predictions distinguishing this from a cosmological constant}

\begin{enumerate}
  \item \textbf{Dark energy is not constant.} It tracks the parent's accretion
    history, which is generically non-constant.
  \item \textbf{Dark energy will eventually reverse.} When the parent enters the
    Hawking-dominated regime (mass loss exceeds accretion), the Cartasis membrane
    shrinks and the baby's expansion decelerates, eventually contracting. This is
    a long-term prediction, not relevant on observational timescales now.
  \item \textbf{Recent dark-energy evolution.} Current observations
    (DESI~\citep{desi2024} and related) hint at evolving dark energy with
    $w_0 w_a$ parameters slightly inconsistent with strict $w = -1$. This is
    potentially consistent with parent-accretion dynamics.
\end{enumerate}

\subsection{Confronting DESI DR2}

The framework makes a sign prediction here, not just a qualitative hope: a parent
that is still accreting but with a \emph{saturating} fractional growth rate gives,
in the CPL parametrization $w(a)=w_0+w_a(1-a)$, a value $w_0>-1$ with $w_a<0$.
\Cref{fig:desi} shows the DESI DR2 fits~\citep{desi2025dr2} against this prediction:
all four dataset combinations (DESI$+$CMB with Pantheon$+$, Union3, or DESY5) lie
in the predicted wedge $w_0>-1,\ w_a<0$, with $\Lambda$CDM at $(-1,0)$ disfavoured
at $2.8$--$4.2\sigma$. The sign pattern the parent-accretion picture predicts is
exactly what is observed---and it follows from the parent black hole growing, with
no cosmological constant.

\paragraph{The phantom crossing, derived.} The DESI fits imply a \emph{phantom}
past ($w<-1$ at $z\gtrsim0.3$) crossing to $w>-1$ today. This is not an extra
assumption to bolt on; it follows from the injection picture directly. Absorb the
membrane injection into an effective equation of state by writing the baby's
dark-energy continuity equation with no explicit source,
$\dot\rho_{\mathrm{DE}}+3H(1+w_{\mathrm{eff}})\rho_{\mathrm{DE}}=0$, which is just the definition
\begin{equation}
  w_{\mathrm{eff}}(a) \;=\; -1 \;-\; \tfrac{1}{3}\,\frac{d\ln\rho_{\mathrm{DE}}}{d\ln a}.
  \label{eq:weff-inj}
\end{equation}
A density that \emph{rises} with $a$ (injection outpacing dilution) gives
$w_{\mathrm{eff}}<-1$; one that \emph{falls} (dilution winning) gives $w_{\mathrm{eff}}>-1$;
a flat one gives exactly $-1$. A parent whose net injection rose and then
saturated therefore makes $\rho_{\mathrm{DE}}(a)$ a hump, and \eqref{eq:weff-inj} forces
$w_{\mathrm{eff}}$ to cross $-1$ \emph{at the peak}. Modelling the bump as a log-normal
in scale factor, $\rho_{\mathrm{DE}}\propto\exp[-\beta(\ln a-\ln a_p)^2]$, gives the
closed form $w_{\mathrm{eff}}(a)=-1+\tfrac{2\beta}{3}(\ln a-\ln a_p)$, which linearises
to CPL with $w_a=-\tfrac{2\beta}{3}$ and $a_p=\exp[(w_0+1)/w_a]$. The DESI$+$CMB$+$
DESY5 values then \emph{predict} the crossing redshift
$z_{\mathrm{cross}}=1/a_p-1\approx0.33$---squarely where DESI's own reconstruction
crosses $-1$. The phantom past is thus a consequence of accretion that rose and
saturated, not a tuned sign. What remains underived is the \emph{amplitude}
($\beta$, equivalently the membrane coupling): we derive the shape and the
crossing, fit the height. The prediction is still falsifiable: if future data
drive $(w_0,w_a)\to(-1,0)$, parent-accretion dark energy fails.

\begin{figure}[t]
  \centering
  \includegraphics[width=\textwidth]{desi_w0wa}
  \caption{Dark energy as parent accretion vs DESI DR2. \emph{Left:} the four
    DESI DR2 $(w_0,w_a)$ fits all fall in the framework's wedge ($w_0>-1$,
    $w_a<0$, shaded); $\Lambda$CDM (star) sits at its corner, disfavoured at
    $2.8$--$4.2\sigma$. \emph{Right:} the peaked-injection $w_{\mathrm{eff}}(z)$ of
    \eqref{eq:weff-inj} crosses $w=-1$ at the injection peak $z_{\mathrm{cross}}\approx
    0.33$ (shaded: phantom epoch), tracking the DESI CPL reconstruction through the
    constrained range. Values flagged for verification against the published
    tables.}
  \label{fig:desi}
\end{figure}

\section{The CMB as parent Hawking radiation}
\label{sec:cmb-hawking}

\subsection{The identification}

The Cartasis membrane is the same geometric surface as the parent's event
horizon (from outside), the baby's cosmic horizon at maximum lookback (from
inside), and the Big Bang ``surface'' in standard cosmology. Hawking radiation
is emitted at the parent's event horizon. From the baby's interior, this same
radiation appears at the cosmic horizon---i.e.\ from every direction at maximum
cosmic lookback distance. \textbf{This matches the observed CMB}: thermal
radiation from every direction at maximum cosmic lookback.

\subsection{Temperature and spectrum}

The Hawking temperature for a black hole of mass $M$ is
\[
  T_H = \frac{\hbar c^3}{8\pi G M k_B}.
\]
For our universe's parent mass ($\sim 10^{53}\,\mathrm{kg}$), the Hawking
temperature is essentially zero at emission. But the radiation has been
redshifting through our cosmic expansion, so the observed temperature is what we
measure: $2.7\,\mathrm{K}$. The redshift factor depends on detailed bounce
geometry; a calculation is needed to verify consistency.

\subsection{Consistency with observations}

The CMB has a near-perfect Planck thermal spectrum (consistent with Hawking
thermality), near-perfect isotropy (consistent with the bounce surface being
homogeneous across the parent's horizon), tiny anisotropies at the $10^{-5}$
level (consistent with bounce-surface inhomogeneities and small parent-rotation
effects), and acoustic peaks in the power spectrum (these need bounce-cosmology
rederivation, not yet done)~\citep{planck2018}.

The standard cosmological explanation of the CMB (recombination at $z \sim 1100$)
fits the spectrum and acoustic features beautifully. The bounce-cosmology
reinterpretation needs to reproduce these features, which is non-trivial.
\textbf{This is the central calculational challenge for the \SCT framework.}

\subsection{Preferred axis}

If our parent is rotating (Kerr-like rather than Schwarzschild-like), the
rotation axis imprints anisotropy on the bounce surface, which manifests as a
preferred direction in the CMB and in subsequent structure formation.
Observational hints of preferred axes include:

\begin{itemize}
  \item the ``axis of evil'' in CMB low multipoles (alignment of the quadrupole
    and octupole with each other and with the ecliptic);
  \item recent JWST galaxy-rotation asymmetry~\citep{shamir2025}: a $\sim 2\!:\!1$
    ratio of clockwise vs.\ counterclockwise rotation in JADES, statistically
    significant at $3.39\sigma$, peaking near the Galactic poles.
\end{itemize}

Both could be signatures of parent rotation. The exact relationship between
parent angular momentum and observed cosmic anisotropies is calculable but not
yet done in detail.

\section{Hubble tension}

The Hubble tension is a persistent $\sim 5\sigma$ disagreement between $H_0$
measured from the CMB ($\sim 67\,\mathrm{km\,s^{-1}\,Mpc^{-1}}$) and from the
local distance ladder ($\sim 73\,\mathrm{km\,s^{-1}\,Mpc^{-1}}$). Standard
cosmology has no clean resolution.

In \SCT, the early universe (near the bounce) is dominated by bounce dynamics
rather than a standard inflation $+$ radiation-dominated phase. The relationship
between early- and late-universe expansion is modified. This could in principle
resolve the Hubble tension, but the specific calculation has not been done. If
the framework can produce the right $H_0$ from late observations while also
matching the CMB precisely, this is a major validation; if it cannot, the
framework needs revision.

\section{Summary of dark-sector predictions}

\begin{table}[ht]
  \centering
  \small
  \begin{tabular}{@{}p{0.18\textwidth}p{0.36\textwidth}p{0.38\textwidth}@{}}
    \toprule
    \textbf{Phenomenon} & \textbf{Standard explanation} & \textbf{\SCT explanation} \\
    \midrule
    Dark matter        & New particle (WIMP, axion, etc.)        & Parent's matter content gravitating through Cartasis \\
    Dark energy        & Cosmological constant $\Lambda$         & Parent-accretion dynamics changing membrane geometry \\
    CMB                & Recombination at $z \sim 1100$          & Parent Hawking radiation redshifted by baby expansion \\
    Matter--antimatter & Sakharov conditions $+$ leptogenesis    & Filter at Cartasis $+$ statistical bias from OG fluctuation \\
    Inflation          & Inflaton field with specific potential  & Torsion-driven expansion at bounce \\
    Hubble tension     & Unresolved                              & Bounce dynamics modify early/late expansion relationship \\
    Preferred axes     & Statistical fluctuation or systematic   & Parent's rotation-axis imprint \\
    \bottomrule
  \end{tabular}
  \caption{Standard vs.\ \SCT interpretations of dark-sector phenomena.}
  \label{tab:dark-sector}
\end{table}

\section{Observational program}

The framework predicts that we have \emph{already observed} the supraverse
through dark-matter and CMB measurements. We just do not recognize these as
supraverse signatures because we interpret them through $\Lambda$CDM.
Reinterpreting through \SCT:

\begin{enumerate}
  \item Dark-matter astronomy is \emph{parent-universe} astronomy. The
    dark-matter distribution encodes information about our parent's structure.
  \item CMB measurements are \emph{parent Hawking radiation} measurements. The
    CMB spectrum and anisotropies encode information about our parent black hole's
    properties.
  \item Dark-energy measurements track \emph{parent accretion history}.
  \item Preferred-axis observations probe \emph{parent rotation}.
\end{enumerate}

This is an enormous research program: re-derive everything we know from
cosmological observations through the \SCT lens, and look for consistencies and
inconsistencies. Much of it is in principle computable from Einstein--Cartan
bounce dynamics applied to specific parent configurations.
